% Physical unit 020; accepted Session 02.
\GordanUnitStart{020}{printed page 15}
\GordanTopFolio{15}
Quo itinere nondum percurso, $\alpha$ et $\beta$ ita possunt determinari, ut functio $N$ finita maneat; excedente vero
\[
 {u\over \sin l_0\sin\varphi}
\]
ultra $0$, alicubi spatii $N$ infinitum exsistat necesse est, ubi illa norma, sub qua
\[
 \int_{\psi_1}^{\psi}d\psi\{(11)\delta l'\delta l'+2(01)\delta l\delta l'+(00)\delta l\delta l\}>0
\]
stare desinit.

Quodsi meminimus,
\[
 {u\over \sin l_0\sin\varphi}
\]
exstitisse, in aequatione lineae geodeticae $\psi$ secundum constantem $l_0$ ita partialiter differentiando, ut $l$ constans intelligatur, ei, quod superiori paragrapho proposuimus, respondet illud theorema.

\begin{quote}
\gq{Linea geodetica tam diu linea brevissima est, quoad aliam proximam, ab eodem puncto exeuntem, lineam geodeticam secet.}
\end{quote}

Punctum hujus sectionis nominatur lineae geodeticae terminus et significatur:
\[
 O(l,\psi).
\]

\begin{center}\S\ 11.\end{center}

Revera
\[
 \int_{\psi_1}^{\psi_2}d\psi\{(11)\delta l'\delta l'+2(01)\delta l\delta l'+(00)\delta l\delta l\}
\]
(posito $\psi=\psi_2$) $\delta l$ ita determinari potest, ut illud evanescat.

Scilicet ponendo
\[
 \delta l=\alpha{\partial l\over \partial l_0}
\]
integrale nostrum formam induit hanc:
\[
 \alpha^2\int_{\psi_1}^{\psi_2}d\psi\left\{(11){\partial l'\over \partial l_0}\cdot{\partial l'\over \partial l_0}
 +2(01){\partial l\over \partial l_0}\cdot{\partial l'\over \partial l_0}
 +(00){\partial l\over \partial l_0}\cdot{\partial l\over \partial l_0}\right\}
\]
vel identitatem:
\[
\begin{aligned}
0={}&\left((10){\partial l\over \partial l_0}{\partial l\over \partial l_0}+(11)'{\partial l\over \partial l_0}{\partial l'\over \partial l_0}\right)
-\left((01){\partial l\over \partial l_0}{\partial l'\over \partial l_0}+(11){\partial l\over \partial l_0}{\partial l'\over \partial l_0}\right),\\
={}&\alpha^2\int_{\psi_1}^{\psi_2}d\psi\left\{-(01)'{\partial l\over \partial l_0}{\partial l\over \partial l_0}
-(11)'{\partial l\over \partial l_0}{\partial l'\over \partial l_0}
-2(01){\partial l\over \partial l_0}{\partial l'\over \partial l_0}
-(11){\partial l'\over \partial l_0}{\partial l'\over \partial l_0}
-(11){\partial l\over \partial l_0}{\partial l''\over \partial l_0}\right\}
\end{aligned}
\]
